\section*{الفصل \ref{ch:b1:topology} --- طوبولوجيا المستقيم الحقيقي}

\begin{solution}{exo:b1:topology:1}
$\intoo{0}{1} \cup \intoo{2}{3}$: \omterm{def:b1:topology:open}{مفتوحة} (اتحاد مجموعات \omterm{def:b1:topology:open}{مفتوحة})، وليست
\omterm{def:b1:topology:closed}{مغلقة} ($\frac 1n \to 0$ خارجها).

$\intco{0}{1}$: لا هذه ولا تلك. ليست \omterm{def:b1:topology:open}{مفتوحة} (فلا \omterm{prop:b1:reals:intervals}{فترة} حول $0$ بالداخل)؛
وليست \omterm{def:b1:topology:closed}{مغلقة} ($1 - \frac1n \to 1 \notin$ \omterm{def:b1:logic:sets}{المجموعة}).

$\{0\} \cup \intcc{1}{2}$: \omterm{def:b1:topology:closed}{مغلقة} (اتحاد منته لمجموعات \omterm{def:b1:topology:closed}{مغلقة})، وليست
\omterm{def:b1:topology:open}{مفتوحة} (إذ تفشل عند $0$).

$\R \setminus \Z$: \omterm{def:b1:topology:open}{مفتوحة} (لأن $\Z$ \omterm{def:b1:topology:closed}{مغلقة})، وليست \omterm{def:b1:topology:closed}{مغلقة}: فالمتتالية
$\bigl(\frac 1n\bigr)$ تقع فيها، لكن نهايتها $0$ تنتمي إلى $\Z$،
أي تفلت من \omterm{def:b1:logic:sets}{المجموعة}.

$\Q \cap \intoo{0}{1}$: لا هذه ولا تلك. ليست \omterm{def:b1:topology:open}{مفتوحة}: فكل \omterm{prop:b1:reals:intervals}{فترة} حول عدد
ناطق تحتوي أعدادًا صماء. وليست \omterm{def:b1:topology:closed}{مغلقة}: فهي تحتوي متتاليات
تؤول إلى العدد الأصمّ $\frac{\sqrt 2}{2}$ (\omterm{def:b1:topology:dense}{بالكثافة}).
\end{solution}

\begin{solution}{exo:b1:topology:2}
$A = \intoc{0}{1} \cup \{2\}$: $\mathring A = \intoo{0}{1}$ و
$\overline A = \intcc{0}{1} \cup \{2\}$ و $\partial A = \{0, 1, 2\}$.

$A = \R \setminus \Q$: $\mathring A = \emptyset$ (لأن كل \omterm{prop:b1:reals:intervals}{فترة}
تحتوي أعدادًا ناطقة)، و $\overline A = \R$ (\omterm{def:b1:topology:dense}{بكثافة}
الأعداد الصماء)، و $\partial A = \R$.

$A = \bigl\{\frac{(-1)^n n}{n+1}\bigr\}$: تؤول الحدود الزوجية إلى $1$،
والفردية إلى $-1$، ولا واحدة منهما تنتمي إلى $A$. ومنه $\mathring A =
\emptyset$ (لأن النقاط معزولة) و $\overline A = A \cup \{-1, 1\}$ و
$\partial A = \overline A$.
\end{solution}

\begin{solution}{exo:b1:topology:3}
\emph{بالمتممات.} $F = \{a_1 < a_2 < \dots < a_k\}$: المتممة
هي اتحاد الفترات \omterm{def:b1:topology:open}{المفتوحة} $\intoo{-\infty}{a_1}$ و
$\intoo{a_i}{a_{i+1}}$ و $\intoo{a_k}{+\infty}$ --- وهي \omterm{def:b1:topology:open}{مفتوحة} حسب
\cref{prop:b1:topology:openstable}.

\emph{بالمتتاليات.} ليكن $u_n \in F$ و $u_n \to \ell$. مع $\varepsilon
= \min\{\abs{a_i - a_j} : i \neq j\}/2 > 0$ (أو مع أيّ $\varepsilon$ إذا كانت
$F$ أحادية): بعد رتبة ما، تقع كل الحدود على بعد
$\varepsilon$ من $\ell$، ومنه على بعد $2\varepsilon$ بعضها من بعض،
وهذا يفرض أن تكون $a_i$ واحدًا ابتداءً من تلك الرتبة؛ وعندئذ
$\ell = a_i \in F$.
\end{solution}

\begin{solution}{exo:b1:topology:4}
$\subseteq$: \omterm{def:b1:logic:sets}{المجموعة} $\overline A \cup \overline B$ \omterm{def:b1:topology:closed}{مغلقة} (اتحاد منته)
وتحتوي $A \cup B$، ومنه فهي تحتوي \emph{أصغر} حاوية \omterm{def:b1:topology:closed}{مغلقة}
$\overline{A \cup B}$. و $\supseteq$: يعطي $A \subseteq A \cup B$
أن $\overline A \subseteq \overline{A \cup B}$ (\omterm{def:b1:topology:closure}{فالغلق}
رتيب: إذ إن النقاط الملاصقة \omterm{def:b1:logic:sets}{للمجموعة} $A$ ملاصقة \omterm{def:b1:logic:sets}{للمجموعة} الأكبر)، و
كذلك من أجل $B$.

ومثال مضاد من أجل التقاطعات: $A = \intoo{0}{1}$ و $B =
\intoo{1}{2}$: $\overline{A \cap B} = \overline\emptyset = \emptyset$
بينما $\overline A \cap \overline B = \{1\}$.
\end{solution}

\begin{solution}{exo:b1:topology:5}
استعمل التمييز التتالي
(\cref{thm:b1:topology:seqclosed}). ليكن $S = \{u_n\} \cup \{\ell\}$
ولتكن $(v_k)$ متتالية في $S$ مع $v_k \to m$؛ نبيّن $m \in S$. وحالتان.
إذا أخذت $(v_k)$ قيمة $v \in S$ ما عددًا لا يُحصى من المرات،
فإن \omterm{def:b1:seq:subsequence}{متتالية جزئية} ثابتة تعطي $m = v \in S$. وإلا فتُؤخذ كل
قيمة عددًا منتهيًا من المرات؛ وعلى الخصوص، من أجل كل $n$، يظهر الحد
$u_n$ عددًا منتهيًا من المرات، وكذلك $\ell$. وعندئذ من أجل كل $N$، تكون
الأدلة $k$ ذات $v_k \in \{u_0, \dots, u_N, \ell\}$ منتهية
العدد: فالحدود $v_k$ الباقية هي حدود $u_n$ مع $n > N$. ومن أجل
$\varepsilon > 0$ معطى، اختر $N$ يحقق $\abs{u_n - \ell} \leq
\varepsilon$ من أجل $n > N$: فتحقق كل الحدود $v_k$ إلا عددًا منتهيًا منها
$\abs{v_k - \ell} \leq \varepsilon$. ومنه $v_k \to \ell$، ومنه $m =
\ell \in S$.
\end{solution}

\begin{solution}{exo:b1:topology:6}
$U + A = \bigcup_{a \in A} (U + a)$، وكل منسحبة $U + a$
\omterm{def:b1:topology:open}{مفتوحة} (بانسحاب شهادات الفترات). واتحاد المجموعات \omterm{def:b1:topology:open}{المفتوحة}
مفتوح (\cref{prop:b1:topology:openstable}).

وأمّا المجموعات \omterm{def:b1:topology:closed}{المغلقة}: فالمجموعتان $\Z$ و $\sqrt 2\,\Z$ مغلقتان (بوصفهما $\alpha\Z$:
لأن المتتاليات المتقاربة ثابتة في النهاية، قارن
\cref{ex:b1:topology:closed}). ومجموعهما $\Z + \sqrt 2\,\Z$ كثيف
في $\R$ (\cref{exo:b1:reals:9}) لكنه ليس $\R$ (فهو قابل للعدّ، أو
ببساطة: $\frac{\sqrt 2}{2} \notin \Z + \sqrt2\Z$، وإلا لكان $\sqrt 2$
ناطقًا --- إذ إن كتابة $\frac{\sqrt2}{2} = m + n\sqrt 2$
تفرض $(2n - 1)\sqrt 2 = -2m$، ومنه $\sqrt 2 \in \Q$ ما لم يكن $n =
\frac12$، وهذا محال). \omterm{def:b1:logic:sets}{والمجموعة} الجزئية القطعية \omterm{def:b1:topology:dense}{الكثيفة} ليست \omterm{def:b1:topology:closed}{مغلقة}: إذ
غلقها هو $\R \neq$ نفسها.
\end{solution}

\begin{solution}{exo:b1:topology:7}
$\Z$: \omterm{def:b1:topology:open}{الجوار} $\intoo{n - \frac12}{n + \frac12}$ للعدد $n$
لا يقاطع $\Z$ إلا عند $n$.

$A = \{\frac1n : n \in \N^*\}$: حول $\frac 1n$، تعزلها عن جارتيها الفترةُ ذات
نصف القطر $\frac{1}{n} - \frac{1}{n+1} = \frac{1}{n(n+1)}$ (منصَّفًا، مثلًا)
--- فكل النقاط معزولة. ومع ذلك $0 \in
\overline A \setminus A$: فالنقاط المعزولة لا تمنع وجود
ملاصقين من الخارج.

وإذا كانت كل نقاط $A$ معزولة: فلا نقطة من $A$ داخلية، لأن
لكل نقطة داخلية \omterm{prop:b1:reals:intervals}{فترة} كاملة من جيران $A$ حولها
(\omterm{prop:b1:reals:intervals}{والفترة} غير منتهية)، وهذا يناقض العزلة. ومنه $\mathring A
= \emptyset$.
\end{solution}

\begin{solution}{exo:b1:topology:8}
إذا كان $A \subseteq \intcc{-M}{M}$، فإن \omterm{def:b1:topology:closed}{المجموعة المغلقة} $\intcc{-M}{M}$
تحتوي $A$، ومنه تحتوي $\overline A$ (وهو أصغر حاوية
\omterm{def:b1:topology:closed}{مغلقة}): ومنه \omterm{def:b1:logic:sets}{فالمجموعة} $\overline A$ محدودة.

وليكن $s = \sup A$ (منتهيًا). ولأن $A \subseteq \overline A$، $\sup
\overline A \geq s$. وبالعكس، $\overline A \subseteq
\intoc{-\infty}{s}$: فنصف المستقيم مغلق ويحتوي $A$؛ ومنه
فكل عنصر من $\overline A$ هو $\leq s$، فينتج $\sup\overline A
\leq s$. ومنه المساواة.

و $s \in \overline A$: فحسب التمييز بالمقدار $\varepsilon$
(\cref{prop:b1:reals:epsilon})، تحتوي كل \omterm{prop:b1:reals:intervals}{فترة} $\intoo{s -
\varepsilon}{s + \varepsilon}$ عنصرًا من $A$: ومنه فالعدد $s$
ملاصق.
\end{solution}

\begin{solution}{exo:b1:topology:9}
افترض أن $A$ \omterm{def:b1:topology:open}{مفتوحة} \omterm{def:b1:topology:closed}{ومغلقة}، مع $a \in A$ و $b \in \R
\setminus A$؛ ودون إخلال بالعمومية $a < b$. \omterm{def:b1:logic:sets}{المجموعة} $B = A \cap
\intcc{a}{b}$ غير خالية (لأن $a$ فيها) ومحدودة: وليكن $s = \sup B$. وحسب
\cref{exo:b1:topology:8}، $s \in \overline B \subseteq \overline A =
A$ (لأن $A$ \omterm{def:b1:topology:closed}{مغلقة}). ولاحظ $s \leq b$، ولأن $b \notin A$: يكون $s < b$.
والآن $A$ \omterm{def:b1:topology:open}{مفتوحة}: فتقع \omterm{prop:b1:reals:intervals}{فترة} ما $\intoo{s - r}{s + r}$ في $A$،
ويمكننا أخذ $r < b - s$. عندئذ ينتمي $s + \frac{r}{2}$ إلى $A
\cap \intcc{a}{b} = B$ ويتجاوز $s$ --- وهذا يناقض $s = \sup
B$. ومنه فلا وجود لمثل هذا الزوج $(a, b)$: فتكون إحدى المجموعتين $A$ و $\R \setminus A$
خالية.
\end{solution}

\begin{solution}{exo:b1:topology:10}
\omterm{def:b1:logic:sets}{المجموعة} $I_x$ اتحاد فترات \omterm{def:b1:topology:open}{مفتوحة} كلها تحتوي $x$: فهي \omterm{def:b1:topology:open}{مفتوحة}،
وهي \omterm{prop:b1:reals:intervals}{فترة} لكونها محدَّبة --- إذ إذا كان $u < z < v$ مع $u, v \in
I_x$، فإن $u$ و $v$ يقعان في فترتين جزئيتين مفتوحتين $J_u \ni x$ و $J_v \ni
x$ من $U$، \omterm{def:b1:logic:sets}{والمجموعة} $J_u \cup J_v$ \omterm{prop:b1:reals:intervals}{فترة} (لأن كلتيهما تحتوي $x$)
داخل $U$ وتحتوي $z$؛ ومنه $z \in I_x$
(\cref{prop:b1:reals:intervals}).

وإذا كان $I_x \cap I_y \neq \emptyset$: فإن $I_x \cup I_y$ عندئذ \omterm{prop:b1:reals:intervals}{فترة}
\omterm{def:b1:topology:open}{مفتوحة} (محدَّبة: فترتان متداخلتان) محتواة في $U$
وتحتوي $x$ و $y$، ومنه $I_x \cup I_y \subseteq I_x$ و
$\subseteq I_y$ بأكبرية كلٍّ منهما: ومنه $I_x = I_y$.

ومنه \omterm{def:b1:logic:sets}{فالمجموعة} $U$ هي الاتحاد المنفصل للمجموعات المتمايزة $I_x$ (فكل $x \in
U$ يقع في $I_x$ الخاصة به). وأمّا قابلية العدّ: فكل \omterm{prop:b1:reals:intervals}{فترة} \omterm{def:b1:topology:open}{مفتوحة} غير خالية
$I$ من العائلة تحتوي عددًا ناطقًا $q_I$
(\cref{thm:b1:reals:density})، وتأخذ الفترات المنفصلة المتمايزة أعدادًا ناطقة
متمايزة: ومنه تنغرس العائلة في $\Q$، وهي قابلة للعدّ
(لأنها مدلَّلة بأزواج من الأعداد الصحيحة). ومنه فعدد الفترات قابل للعدّ
على الأكثر.
\end{solution}

\begin{solution}{exo:b1:topology:11}
$\overline A = A \cup A'$. ($\supseteq$) $A \subseteq \overline
A$ دائمًا؛ وإذا كان $x \in A'$، فإن كل \omterm{def:b1:topology:open}{جوار} للنقطة $x$ يقاطع $A
\setminus \{x\} \subseteq A$، ومنه فالنقطة $x$ ملاصقة. ($\subseteq$)
ليكن $x \in \overline A$. فإذا كان $x \in A$، فقد انتهينا. وإذا كان $x \notin A$،
فإن كل \omterm{def:b1:topology:open}{جوار} للنقطة $x$ يقاطع $A = A \setminus \{x\}$: $x \in
A'$.

ونتيجةً لذلك تكون $A$ \omterm{def:b1:topology:closed}{مغلقة} $\iff$ $A = \overline A = A \cup A'$
$\iff$ $A' \subseteq A$.

$A = \{\frac 1n\}$: العدد $0$ نقطة تراكم ($\frac 1n \to
0$، والحدود $\neq 0$)؛ وكل $\frac 1n$ معزولة
(\cref{exo:b1:topology:7})، ومنه ليست في $A'$؛ ولأيّ نقطة $x
\notin A \cup \{0\}$ \omterm{prop:b1:reals:intervals}{فترة} كاملة تتجنّب $A$ (بين
جارتَي $x$ في $A \cup \{0\}$، أو بعد $1$).
ومنه $A' = \{0\}$.

$\Z' = \emptyset$: فكل عدد صحيح معزول، ولكل عدد غير صحيح
\omterm{def:b1:topology:open}{جوار} داخل $\R \setminus \Z$.

$\Q' = \R$: فكل \omterm{prop:b1:reals:intervals}{فترة} حول أيّ عدد حقيقي تحتوي أعدادًا ناطقة
لا تُحصى (\cref{thm:b1:reals:density})، وعلى الخصوص واحدًا
مختلفًا عن المركز.
\end{solution}

\begin{solution}{exo:b1:topology:12}
\begin{enumerate}
  \item من أجل كل $a \in A$: $\abs{x - a} \leq \abs{x - y} +
        \abs{y - a}$، ومنه $d_A(x) \leq \abs{x - y} + \abs{y - a}$؛
        وبأخذ \omterm{def:b1:reals:bounds}{الحد الأدنى} على $a$: $d_A(x) \leq \abs{x - y} +
        d_A(y)$. وبتبادل $x$ و $y$ نجد المتراجحة
        الأخرى: $\abs{d_A(x) - d_A(y)} \leq \abs{x - y}$.
  \item $d_A(x) = 0$ $\iff$ من أجل كل $\varepsilon > 0$ يوجد
        $a \in A$ يحقق $\abs{x - a} < \varepsilon$ $\iff$ تقاطع كل
        \omterm{prop:b1:reals:intervals}{فترة} حول $x$ المجموعةَ $A$ $\iff$ $x \in \overline A$.
        وإذا كانت $F$ \omterm{def:b1:topology:closed}{مغلقة} و $x \notin F = \overline F$، فإن
        $d_F(x) \neq 0$، أي $d_F(x) > 0$.
  \item إذا كان $d_F(x) < \varepsilon$، فضع $r = \varepsilon -
        d_F(x) > 0$: فمن أجل $\abs{y - x} < r$، يعطي البند (1) أن
        $d_F(y) \leq d_F(x) + \abs{x - y} < \varepsilon$: ومنه تقع
        \omterm{prop:b1:reals:intervals}{الفترة} $\intoo{x - r}{x + r}$ في $V_\varepsilon$،
        وهي إذن \omterm{def:b1:topology:open}{مفتوحة}؛ وهي تحتوي $F$ لأن $d_F = 0$
        هناك. وأخيرًا $x \in \bigcap_{\varepsilon>0}
        V_\varepsilon$ $\iff$ $d_F(x) < \varepsilon$ من أجل كل
        $\varepsilon$ $\iff$ $d_F(x) = 0$ $\iff$ $x \in \overline
        F = F$. ولأن $\bigcap_{\varepsilon > 0} V_\varepsilon =
        \bigcap_{n \geq 1} V_{1/n}$، تكون كل \omterm{def:b1:topology:closed}{مجموعة مغلقة}
        تقاطعًا قابلًا للعدّ لمجموعات \omterm{def:b1:topology:open}{مفتوحة}.
\end{enumerate}
\end{solution}

\begin{solution}{pb:b1:topology:1}
\textbf{1.} $C_1 = \intcc{0}{\frac13} \cup \intcc{\frac23}{1}$
و
\[
C_2 = \intcc{0}{\tfrac19} \cup \intcc{\tfrac29}{\tfrac13}
\cup \intcc{\tfrac23}{\tfrac79} \cup \intcc{\tfrac89}{1} .
\]
وبالاستقراء: إذا كانت $C_n$ اتحادًا منفصلًا لعدد $2^n$ من القطع \omterm{def:b1:topology:closed}{المغلقة}
طول كلٍّ منها $3^{-n}$، فإن حذف الثلث الأوسط المفتوح من كل قطعة
يترك قطعتين مغلقتين طول كلٍّ منهما $3^{-n-1}$ لكل أب:
أي $2^{n+1}$ قطعة، منفصلة مثنى مثنى (فأبناء آباء متمايزين
منفصلون لأن الآباء كانوا كذلك؛ وابنا أب
واحد تفصل بينهما الفجوة المحذوفة).

\textbf{2.} كل $C_n$ اتحاد منته لقطع، ومنه فهي
\omterm{def:b1:topology:closed}{مغلقة}؛ و $C = \bigcap C_n$ تقاطع مجموعات \omterm{def:b1:topology:closed}{مغلقة}:
فهي \omterm{def:b1:topology:closed}{مغلقة} (\cref{def:b1:topology:closed})؛ ومحدودة ($\subseteq
\intcc{0}{1}$): ومنه متراصة حسب \cref{thm:b1:topology:compact}.
وغير خالية: إذ يقع $0$ في القطعة اليسرى القصوى من كل $C_n$. وليكن
$a$ طرفًا لقطعة $S$ من $C_n$. فمن أجل $m \leq n$،
$a \in C_n \subseteq C_m$. وأمّا المراحل اللاحقة: فحذف الثلث
الأوسط لا يزيل طرفًا أبدًا، ويكون $a$ مرة أخرى طرفًا
لأحد ابنَي $S$ (الابن الملامس للنقطة $a$)؛ ومنه
بالاستقراء $a \in C_m$ من أجل كل $m \geq n$: أي $a \in C$.

\textbf{3.} الطول الكلي للمقدار $C_n$: $2^n \cdot 3^{-n} =
(\frac23)^n \to 0$. ومن أجل $\varepsilon > 0$ معطى، اختر $n$ يحقق
$(\frac23)^n \leq \varepsilon$: عندئذ $C \subseteq C_n$، وهي اتحاد
عدد منته من القطع طولها الكلي $\leq \varepsilon$.

\textbf{4.} لتكن $I \subseteq C$ \omterm{prop:b1:reals:intervals}{فترة} فيها نقطتان
متمايزتان. من أجل كل $n$: $I \subseteq C_n$، وعلى $I$، لكونها محدَّبة،
أن تقع داخل قطعة \emph{واحدة} من $C_n$ (إذ إن مقاطعة قطعتين
ستفرض على $I$ أن تحتوي نقطة من الفجوة بينهما،
وهي خارج $C_n$). ومنه فطول $I$ هو $\leq
3^{-n}$ من أجل كل $n$: وهذا تناقض. ومنه فالفترات الوحيدة
داخل $C$ خالية أو أحادية؛ وعلى الخصوص لا
$\intoo{x-r}{x+r}$ تسع داخل $C$: $\mathring C = \emptyset$.
ولأن $C$ \omterm{def:b1:topology:closed}{مغلقة}، يكون داخل $\overline C = C$ خاليًا: ومنه \omterm{def:b1:logic:sets}{فالمجموعة} $C$
غير \omterm{def:b1:topology:dense}{كثيفة} في أيّ موضع.

\textbf{5.} اكتب $\varphi_0(x) = \frac x3$ و $\varphi_2(x) =
\frac{2 + x}{3}$، وهما تقابلان أفينيان متزايدان من $\intcc{0}{1}$
على $\intcc{0}{\frac13}$ و $\intcc{\frac23}{1}$. والادعاء:
$C_{n+1} = \varphi_0(C_n) \cup \varphi_2(C_n)$. ومن أجل $n = 0$ هذا
هو السؤال 1. وبالاستقراء: يرسل \omterm{def:b1:logic:map}{تطبيق} أفيني متزايد
الثلثَ الأوسط لقطعة إلى الثلث الأوسط للقطعة
الصورة، ومنه فحذف الأثلاث الوسطى يتبادل مع $\varphi_0$ و
$\varphi_2$؛ \omterm{def:b1:logic:map}{وتطبيق} خطوة الحذف على $C_{n+1} =
\varphi_0(C_n) \cup \varphi_2(C_n)$ يعطي $C_{n+2} =
\varphi_0(C_{n+1}) \cup \varphi_2(C_{n+1})$. وبالتقاطع على
$n$: من أجل $x \leq \frac13$، $x \in C \iff x \in \varphi_0(C_n)$
من أجل كل $n$ $\iff 3x \in \bigcap C_n = C$؛ وكذلك على
$\intcc{\frac23}{1}$؛ ولا نقطة من $\intoo{\frac13}{\frac23}$
تقع في $C_1$. ومنه $C = \varphi_0(C) \cup \varphi_2(C)$،
انفصالًا.

\textbf{6.} بالاستقراء على $n$؛ وحالة $n = 0$ تقول إن لكل $x \in
\intcc{0}{1}$ ترميزًا، وهو \cref{pb:b1:reals:1}
(السؤال 9 من أجل $x < 1$؛ $1 = (0.\overline 2)_3$). افترض
التكافؤ عند الرتبة $n$. فإذا كان $x \in C_{n+1}$: فحسب السؤال 5، $x =
\varphi_i(z)$ مع $z \in C_n$ و $i \in \{0, 2\}$؛ وإذا كان $(e_k)$
ترميزًا للمقدار $z$ أرقامه $n$ الأولى خالية من $1$، فإن $(i, e_1,
e_2, \dots)$ له مجاميع جزئية $\frac i3 + \frac13\sum_{k\leq m}
e_k 3^{-k} \to \varphi_i(z) = x$: أي ترميز للمقدار $x$ أرقامه $n +
1$ الأولى خالية من $1$. وبالعكس، إذا كان للمقدار $x$ ترميز $(d_k)$ مع
$d_1, \dots, d_{n+1} \in \{0, 2\}$: فإن السلسلة المزاحة $(d_2,
d_3, \dots)$ لها قيمة $z \in \intcc{0}{1}$، وأرقامها $n$
الأولى خالية من $1$، ويعطي حساب المجاميع الجزئية مقروءًا
بالمقلوب أن $x = \varphi_{d_1}(z)$؛ وبالاستقراء $z \in
C_n$، ومنه $x \in C_{n+1}$ حسب السؤال 5. وأخيرًا: يضع ترميز
خالٍ من $1$ كليًا المقدارَ $x$ في كل $C_n$، ومنه في $C$؛ وبالعكس
إذا كان $x \in C$، فإنه من أجل كل $n$ يكون لأحد ترميزَي
$x$ على الأكثر (\cref{pb:b1:reals:1}، السؤال 11) أرقامُه $n$ الأولى
خالية من $1$؛ ولا بد أن يفي ترميز مثبَّت واحد بالغرض من أجل $n$ كبير
كما نشاء (بمبدأ الأدراج بين ترميزين)، والترميز الذي أرقامه $n$
الأولى خالية من $1$ من أجل $n$ كبير كما نشاء يكون خاليًا من $1$
تمامًا.

\textbf{7.} $0 = (0.\overline 0)_3$ و $1 = (0.\overline 2)_3$ و
$\frac13 = (0.0\overline{2})_3$ (وهو التوأم غير السليم للمقدار $(0.1)_3$) و
$\frac23 = (0.2\overline{0})_3$. وبالمجاميع الهندسية:
\[
(0.\overline{02})_3 = \sum_{j\geq1} \frac{2}{9^{\,j}}
= \frac{2/9}{1 - 1/9} = \frac14 ,
\qquad
(0.\overline{20})_3 = \sum_{j\geq1} \frac{2}{3\cdot 9^{\,j-1}}
= \frac{2/3}{1 - 1/9} = \frac34 ,
\]
وكلاهما خالٍ من $1$: $\frac14, \frac34 \in C$. (والمجاميع غير المنتهية
اختصارٌ للحدود العليا للمجاميع الجزئية، كما في \cref{pb:b1:reals:1}.)
وأطراف قطع $C_n$ لها الشكل $m/3^n$
(بالاستقراء: فأطراف الأبناء أطرافُ آباء أو تختلف
عن واحد منها بمضاعف للمقدار $3^{-n-1}$). وإذا كان $\frac14 =
\frac{m}{3^n}$ فإن $3^n = 4m$، و $4 \nmid 3^n$: وهذا محال.
ومنه $\frac14 \in C$ دون أن يكون طرفًا أبدًا.

\textbf{8.} افترض أن للمقدار $x$ ترميزين متمايزين خاليين من $1$. وكونُ
له ترميزان أصلًا يعني (\cref{pb:b1:reals:1}، السؤال 11، الأساس
$3$) أن $x = m/3^N \in \intoo{0}{1}$ وأن الترميزين هما:
المنتهي، وآخر رقم غير معدوم فيه $d_N \in \{1, 2\}$
متبوعًا بأرقام $0$، وتوأمه، وفيه $d_N - 1$ في الموضع $N$
متبوعًا بأرقام $2$. فإذا كان $d_N = 1$ احتوى الأول رقم $1$؛ وإذا كان $d_N
= 2$ حمل التوأم $d_N - 1 = 1$. وفي الحالتين لا يكون خاليًا من $1$
إلا واحد من الزوج على الأكثر: وهذا تناقض. ومنه فلكل $x \in C$
ترميز واحد بالضبط خالٍ من $1$ (والوجود من السؤال 6)، و
لسلاسل $\{0,2\}$ المتمايزة قيمٌ متمايزة. ولكل
سلسلة $\{0,2\}$ قيمةٌ في $\intcc{0}{1}$ (فالمجاميع الجزئية $\leq
1$) وكل بادئاتها خالية من $1$، ومنه فقيمتها في كل $C_n$،
أي في $C$: ومنه \omterm{def:b1:logic:map}{فتطبيق} القيمة تقابلٌ من
سلاسل $\{0,2\}$ على $C$.

\textbf{9.} ليكن $(d^{(k)})$ الترميزَ الخالي من $1$ للمقدار $x_k$ و
ضع $e_k = 2 - d^{(k)}_k \in \{0, 2\}$: وهي سلسلة $\{0,2\}$
قيمتها $y$ تقع في $C$ ولها $(e_k)$ ترميزًا وحيدًا خاليًا من
$1$ (السؤال 8). ومن أجل كل $k$ يختلف ترميزا $y$ و $x_k$
عند الموضع $k$، ومنه $y \neq x_k$: فلا \omterm{def:b1:logic:map}{تطبيق} $\N^* \to C$
شامل. أي \omterm{def:b1:logic:sets}{مجموعة} غير قابلة للعدّ وطولها معدوم: كبرٌ في
القوة العددية وصغرٌ في القياس --- في آن واحد.

\textbf{10.} التراص: انغلاق التقاطع غير المنتهي
مع الحدّية (السؤال 2) --- وهي الخاصية الوحيدة
التي لا يحملها احتواء واحد. والطول المعدوم: $C \subseteq
C_n$ وطولها الكلي $(\frac23)^n$ (السؤال 3). وأمّا عدم
\omterm{def:b1:topology:dense}{الكثافة} في أيّ موضع: فإن $C \subseteq C_n$ يفرض على الفترات داخل $C$ أن يكون
طولها $\leq 3^{-n}$ (السؤال 4). وأمّا عدم قابلية العدّ فلا تركب على أيّ
احتواء أصلًا: إذ تحتاج إلى بنية التقاطع الكاملة،
مرمَّزةً في تقابل السؤال 8.

\textbf{11.} اقلب $d_n$ إلى $2 - d_n$: فتبقى السلسلة الجديدة
سلسلةَ $\{0,2\}$، ومنه تقع قيمتها $x_n$ في $C$؛ والمجاميع
الجزئية بعد الرتبة $n$ تختلف بالمقدار $2\cdot3^{-n}$ بالضبط، ومنه
$\abs{x_n - x} = 2\cdot3^{-n}$. ومنه $x_n \neq x$ و $x_n \to
x$: فكل نقطة من $C$ نهايةٌ لنقاط أخرى من $C$ --- ومنه \omterm{def:b1:logic:sets}{فالمجموعة} $C$
\emph{تامّة}، بلا نقطة معزولة.

\textbf{12.} حسب استقراء السؤال 5، تكون قطع
$C_n$ هي بالضبط $\intcc{t}{t + 3^{-n}}$ حيث يجري $t$
على قيم سلاسل $\{0,2\}$ ذات الطول $n$. ومن أجل $x \in C$
بالترميز $(d_k)$، يكون الاقتطاع $t_n$ (الأرقام $d_1 \dots d_n$
ثم أرقام $0$) طرفًا أيسر إذن، و $0 \leq x - t_n
\leq 3^{-n}$: فالأطراف \omterm{def:b1:topology:dense}{كثيفة} في $C$. وهي تكوّن \omterm{def:b1:logic:sets}{مجموعة} جزئية من
$\{m/3^n : m, n\}$، وهي \omterm{def:b1:logic:sets}{مجموعة} مدلَّلة بأزواج من الأعداد الصحيحة، ومنه
قابلة للعدّ (كما من أجل $\Q$ في \cref{exo:b1:topology:10}). ولأن $C$
غير قابلة للعدّ (السؤال 9)، فإن كل نقاط $C$ إلا عددًا قابلًا للعدّ منها
ليست أطرافًا --- والعدد $\frac14$ (السؤال 7) هو القمة المرئية
من ذلك الجبل الجليدي.

\textbf{13.} اختر $n$ يحقق $3^{-n} < y - x$. فكلاهما $x, y \in
C_n$، ولا يمكن أن يقعا في القطعة نفسها (لأن طولها $3^{-n} <
y - x$): ومنه توفّر الفجوة المحذوفة بين قطعتيهما نقطةً $z$
تحقق $x < z < y$ و $z \notin C_n \supseteq C$. ومنه فأيّ نقطتين
من $C$ تفصل بينهما المتممة: ومنه فالمجموعات الجزئية المحدَّبة الوحيدة
من $C$ أحادية --- أي إن $C$ مفكّكة كليًا.

\textbf{14.} إذا كان $(d_k)$ هو الترميز الخالي من $1$ للمقدار $x$، فإن السلسلة
$(2 - d_k)$ سلسلة $\{0,2\}$ مرة أخرى، ومجاميعها الجزئية
\[
\sum_{k=1}^{n} (2 - d_k)3^{-k} = (1 - 3^{-n}) - \sum_{k=1}^n
d_k 3^{-k} \longrightarrow 1 - x :
\]
ومنه $1 - x \in C$. ومنه $1 - C \subseteq C$، وبتطبيق \omterm{def:b1:logic:map}{التطبيق}
مرتين نجد $1 - C = C$: \omterm{pb:b1:topology:1}{فمجموعة كانتور} متناظرة بالنسبة إلى
$\frac12$.

\textbf{15.} المجاميع الجزئية $t_n \to x$ و $t'_n \to x'$
(لأن المتتاليات المتزايدة تتقارب إلى حدّها الأعلى، أي إلى
القيمة)، ومنه $t_n + t'_n \to x + x'$ حسب
\cref{thm:b1:seq:operations}. ومن أجل $y \in \intcc{0}{1}$ بالترميز
$(e_k)$، اختر $(a_k, b_k) = (0,0), (0,2), (2,2)$
تبعًا لكون $e_k = 0, 1, 2$: عندئذ $a_k + b_k = 2e_k$، وتكون
السلسلتان $(a_k)$ و $(b_k)$ سلسلتَي $\{0,2\}$ قيمتاهما $x,
x' \in C$، و
\[
x + x' = \lim_n\,(t_n + t'_n) = \lim_n 2\sum_{k=1}^n e_k 3^{-k}
= 2y :
\]
فيكون كل $y \in \intcc{0}{1}$ منتصف نقطتين من $C$.

\textbf{16.} يعطي السؤال 15 أن $\intcc{0}{2} = 2\,\intcc{0}{1}
\subseteq C + C$، و $C + C \subseteq \intcc{0}{1} +
\intcc{0}{1} = \intcc{0}{2}$: ومنه المساواة. ثم، باستعمال $1 - C = C$:
\[
C - C = C + (C - 1) = (C + C) - 1 = \intcc{-1}{1} .
\]
ومثال ملموس: $1 = \frac14 + \frac34$، وهو مجموع عضوين
غير طرفيين من $C$.

\textbf{17.} يقيس الطول مقدار ما تشغله \omterm{def:b1:logic:sets}{المجموعة} نفسها
من المستقيم؛ وهو لا يقول شيئًا عن \omterm{def:b1:logic:sets}{مجموعة} \emph{المجاميع}، وهي
صورة العائلة ذات الوسيطين $C \times C$ \omterm{def:b1:logic:map}{بالتطبيق} $(x,
x') \mapsto x + x'$ --- فسلسلتا الأرقام تُختاران
مستقلتين، وتلك الحرية هي بالضبط ما يملأ
$\intcc{0}{2}$. ولا مبرهنة تحدّ طول \omterm{def:b1:logic:sets}{مجموعة} المجاميع بأطوال
حدودها، \omterm{def:b1:logic:sets}{والمجموعة} $C$ برهانٌ على أن لا مبرهنة كهذه ممكنة.

\textbf{18.} حسب \cref{pb:b1:reals:1} (السؤال 18)، يكون $x$
ناطقًا إذا وفقط إذا كان نشره السليم دوريًا في النهاية. و
الترميز الخالي من $1$ للمقدار $x \in C$ إمّا ذلك النشر السليم وإمّا
التوأم غير السليم لنشر منته؛ والسلسلة المنتهية و
توأمها (المساوي $2$ في النهاية) كلاهما دوري في النهاية،
ومنه فدورية الترميز الخالي من $1$ تكافئ
نطقية $x$. والقسمة المطوّلة للمقدار $\frac1{13}$ في الأساس $3$
($r_0 = 1$): $3 = 13\cdot0 + 3$ و $9 = 13\cdot0 + 9$ و $27 =
13\cdot2 + 1$، ويعود الباقي إلى $1$: فالأرقام
$\overline{002}$، ومنه $\frac1{13} = (0.\overline{002})_3$،
وهو خالٍ من $1$ ودوري: أي عضو ناطق في $C$. (وللتحقق:
$\frac{2/27}{1 - 1/27} = \frac{2}{26} = \frac1{13}$.)

\textbf{19.} السلسلة ذات $d_k = 2$ في المواضع
المثلثية $k = \frac{j(j+1)}{2}$ و $0$ فيما عداها هي
سلسلة $\{0,2\}$، ومنه فقيمتها $x^*$ تنتمي إلى $C$ (السؤال
8). وفيها أرقام $2$ لا تُحصى بفجوات $j + 1 \to \infty$
بين المتتالية منها، ومنه فهي ليست دورية في النهاية (إذ إن
دورًا $T$ سيفرض في النهاية أرقام $2$ بفجوات $\leq T$: وهي حجة
تنامي الفجوات في \cref{pb:b1:reals:1}، السؤال 20)؛ وحسب
السؤال 18، $x^* \notin \Q$. وحسب السؤال 9 مع
قابلية عدّ $\Q$، تكون كل أعضاء $C$ إلا عددًا قابلًا للعدّ منها
صمّاء: فالمقدار $x^*$ هو القاعدة لا الاستثناء.

\textbf{20.} ليكن $y \in \intcc{0}{1}$: فله ترميز ثنائي
$(c_k)$ مع $c_k \in \{0, 1\}$ (\cref{pb:b1:reals:1}، السؤال
9، الأساس $2$؛ ويأخذ $y = 1$ السلسلة التي كل أرقامها $1$). عندئذ $(2c_k)$
سلسلة $\{0,2\}$، وقيمتها $x$ تقع في $C$، ويكون $h(x)$
قيمةَ $(c_k)$، أي $y$: ومنه يرسل $h$ المجموعةَ $C$ على
$\intcc{0}{1}$. ولو كانت $C$ صورة \omterm{def:b1:logic:map}{تطبيق} من $\N^*$،
لكان التركيب مع $h$ يعدّد كل $\intco{0}{1}$،
مناقضًا مبرهنة القطر في \cref{pb:b1:reals:1}
(السؤال 22): ومنه \omterm{def:b1:logic:sets}{فالمجموعة} $C$ غير قابلة للعدّ، مرة أخرى. أي \omterm{def:b1:logic:sets}{مجموعة} طولها معدوم
تشمل قطعة كاملة.

\textbf{21.} حسب السؤال 5، تكون $C$ الاتحادَ المنفصل للمجموعتين
$\varphi_0(C)$ و $\varphi_2(C)$، وكلٌّ منهما منسحبة النسخة
المصغَّرة $\frac13 C$. وتعطي الجمعية والسلّمة والحفظ
بالانسحاب أن
\[
L(C) = L(\varphi_0(C)) + L(\varphi_2(C))
= \tfrac13 L(C) + \tfrac13 L(C) = \tfrac23\,L(C),
\]
ومنه $\frac13 L(C) = 0$: أي $L(C) = 0$. فالتشابه الذاتي وحده
يحكم على $C$ بطول معدوم --- ولم يفعل السؤال 3 إلا أن نفّذ
الحكم.

\textbf{22.} تفقد قطعة طولها $l_n$ \omterm{prop:b1:reals:intervals}{فترة} مركزية
طولها $4^{-(n+1)}$، فتبقى قطعتان طول كلٍّ منهما $l_{n+1}
= \frac{l_n - 4^{-(n+1)}}{2}$؛ وانطلاقًا من $l_0 = 1$، يؤكّد الاستقراء
$l_n = \frac{2^n + 1}{2\cdot4^n}$: فعلًا
$\frac12\Bigl(\frac{2^n+1}{2\cdot4^n} - \frac{1}{4^{n+1}}\Bigr)
= \frac{2(2^n + 1) - 1}{2\cdot4^{n+1}} = \frac{2^{n+1} +
1}{2\cdot4^{n+1}}$، و $l_n > 0$ دائمًا: فالبناء لا
يجوع أبدًا. \omterm{def:b1:logic:sets}{والمجموعة} $K = \bigcap K_n$ \omterm{def:b1:topology:closed}{مغلقة} ومحدودة، ومنه متراصة؛
وكل \omterm{prop:b1:reals:intervals}{فترة} داخل $K$ تقع في قطعة واحدة من $K_n$، طولها
$l_n \to 0$: فداخلها خالٍ. والطول المحذوف: $\sum_{n\geq0} 2^n
\cdot 4^{-(n+1)} = \frac14\sum_{n\geq0}\bigl(\frac12\bigr)^n =
\frac12$، ولكل $K_n$ طول كلي $2^n l_n = \frac{2^n +
1}{2^{n+1}} > \frac12$. والآن ليكن عدد منته من الفترات \omterm{def:b1:topology:open}{المفتوحة}
اتحادها $U \supseteq K$. فحسب \omterm{def:b1:logic:statement}{العبارة} الرفيقة في
\cref{ex:b1:topology:nested}، $U \supseteq K_n$ من أجل $n$ ما؛
وبقبول جمعية الطول على الاتحادات المنتهية للفترات،
يكون الطول الكلي لفترات التغطية على الأقل طولَ
$K_n$، وهو يتجاوز $\frac12$. ومنه \omterm{def:b1:logic:sets}{فالمجموعة} $K$ غير \omterm{def:b1:topology:dense}{كثيفة} في أيّ موضع، ومع ذلك لا
تغطية رخيصة لها: فالصغر الطوبولوجي (غير \omterm{def:b1:topology:dense}{كثيفة} في أيّ موضع) و
الصغر المتري (طول معدوم) مفهومان مختلفان حقًا،
\omterm{def:b1:logic:sets}{والمجموعة} $K$ تفصل بينهما.

\textbf{23.} البلوغ: ليكن $d = d(x, F)$ واختر $a_k \in F$
يحقق $\abs{x - a_k} \leq d + \frac1k$: فالمقادير $a_k$ محدودة، ومنه
تستخرج بولتزانو--فايرشتراس (\cref{thm:b1:seq:bw})
$a_{\varphi(k)} \to a$، مع $a \in F$ (لأن $F$ \omterm{def:b1:topology:closed}{مغلقة}،
\cref{thm:b1:topology:seqclosed}) و $\abs{x - a} = \lim
\abs{x - a_{\varphi(k)}} = d$. وأمّا القيمة العظمى: فإذا كان $y \in C$،
فإن $d(y, C) = 0$؛ وإلا وقع $y$ في فجوة محذوفة عند مرحلة ما
$n \geq 1$، وهي \omterm{prop:b1:reals:intervals}{فترة} \omterm{def:b1:topology:open}{مفتوحة} طولها $3^{-n}$ وطرفاها
ينتميان إلى $C$ (السؤال 2)، ومنه $d(y, C) \leq
\frac{3^{-n}}{2} \leq \frac16$، والمساواة تقتضي $n = 1$
و $y$ في مركز الفجوة $\intoo{\frac13}{\frac23}$،
أي $y = \frac12$؛ وفعلًا $d\bigl(\frac12, C\bigr) =
\frac16$ لأن $C \cap \intoo{\frac13}{\frac23} = \emptyset$ و
$\frac13, \frac23 \in C$. ومنه $\max_{y\in\intcc{0}{1}} d(y, C)
= \frac16$، مبلوغًا بالضبط عند $\frac12$.

\textbf{24.} تكوّن الأطراف \omterm{def:b1:logic:sets}{مجموعة} جزئية قابلة للعدّ \omterm{def:b1:topology:dense}{وكثيفة} في $C$
(السؤال 12): عدّدها متتاليةً واحدة $(e_j)_{j\geq1}$،
وهي متتالية في $C$. ونهايات متتالياتها الجزئية كلها تقع في $C$ (لأن $C$
\omterm{def:b1:topology:closed}{مغلقة}). وبالعكس، ثبّت $x \in C$: فمن أجل كل $n$، تقع أطراف قطع
المجموعات $C_m$ التي تحتوي $x$ ($m \geq n$) على بعد
$3^{-m} \leq 3^{-n}$ من $x$، ومنه يقع عدد لا يُحصى من الأطراف
المتمايزة على بعد $3^{-n}$ من $x$؛ فاختر أدلة $j_1 < j_2
< \dots$ تحقق $\abs{e_{j_n} - x} \leq 3^{-n}$: فتكون \omterm{def:b1:seq:subsequence}{متتالية جزئية}
تتقارب إلى $x$. ومنه \omterm{def:b1:logic:sets}{فمجموعة} نهايات المتتاليات الجزئية للمتتالية $(e_j)$
هي $C$ بالضبط --- متتالية واحدة تتكدّس عند نقاط لا تُحصى
عددًا غير قابل للعدّ، وهو الطرف المعاكس لمتتالية متقاربة
\omterm{def:b1:logic:sets}{مجموعة} تكدّسها أحادية.

\textbf{25.} (أ) استهلك البناء: استقرار المجموعات
\omterm{def:b1:topology:closed}{المغلقة} بالتقاطع الكيفي (لوجود $C$ مجموعةً
\omterm{def:b1:topology:closed}{مغلقة})، ومبرهنة التراص \cref{thm:b1:topology:compact}
(الأسئلة 2 و 22 و 23)، والتمييزات التتالية
للانغلاق وللملاصقة (حجة المتراصات المتداخلة في
\cref{ex:b1:topology:nested} والسؤال 23). (ب) والأزواج
الأربعة: طول معدوم ومع ذلك غير قابلة للعدّ (السؤالان 3 و 9)؛ \omterm{def:b1:topology:closed}{ومغلقة}
ومع ذلك داخلها خالٍ (السؤال 4)؛ وتامّة --- بلا نقطة
معزولة --- ومع ذلك مفكّكة كليًا (السؤالان 11 و 13)؛
ومهملة ومع ذلك $C + C = \intcc{0}{2}$ (السؤال 16). (ج)
\omterm{def:b1:logic:inj}{والتطبيق الشامل} $h$ من السؤال 20، مجعولًا متصلًا و
غير متناقص، يصير سلّم الشيطان في نظرية
الدوال المتصلة؛ وفي نظرية القياس في مجلّد السنة الثالثة،
تكون $C$ الشاهدَ المعياري على أن «مهملة» لا تعني
«قابلة للعدّ»، مع ابن عمّها السمين (السؤال 22) الذي يفصل
بين «غير \omterm{def:b1:topology:dense}{كثيفة} في أيّ موضع» و«مهملة».
\end{solution}
